% generated by npm run cv:generate; edit src/content, not this file
\documentclass[10pt,a4paper]{article}
\usepackage[a4paper,left=1.5cm,right=1.5cm,top=1.1cm,bottom=1.1cm]{geometry}
\usepackage[T1]{fontenc}
\usepackage[utf8]{inputenc}
\usepackage[scaled=0.92]{helvet}
\renewcommand{\familydefault}{\sfdefault}
\usepackage{xcolor}
\usepackage{enumitem}
\usepackage{titlesec}
\usepackage{needspace}
\usepackage[hidelinks]{hyperref}

\definecolor{ink}{HTML}{222222}
\definecolor{muted}{HTML}{555555}
\definecolor{rule}{HTML}{BBBBBB}
\color{ink}
\pagestyle{empty}
\setlength{\parindent}{0pt}
\setlength{\parskip}{0pt}
\raggedright

\titleformat{\section}{\normalsize\bfseries}{}{0pt}{\MakeUppercase}[{\color{rule}\titlerule[0.4pt]}]
\titlespacing*{\section}{0pt}{7pt}{3pt}

\setlist[itemize]{leftmargin=1.1em,label=\textbullet,itemsep=0pt,topsep=1pt,parsep=0pt,partopsep=0pt}

% \role{title}{organisation}{place and dates}
\newcommand{\role}[3]{\vspace{5pt}\needspace{3\baselineskip}\textbf{#1}\,\,|\,\,#2\hfill{\small\color{muted}#3}\par}
% \rolenote{text}
\newcommand{\rolenote}[1]{{\small\itshape\color{muted}#1}\par}
% \project{name}{context}
\newcommand{\project}[2]{\vspace{4pt}\needspace{3\baselineskip}\textbf{#1}\hfill{\small\color{muted}#2}\par}
% \stack{tools}
\newcommand{\stack}[1]{{\small\itshape\color{muted}\hspace{1.1em}\textbf{Stack:} #1}\par}

% \cvphone: set by the PDF build from cv/phone.txt; empty when unset
\providecommand{\cvphone}{}

\begin{document}

{\LARGE\bfseries KAREEM FAKHRY GAHEED ABD-ALLAH}\par\vspace{3pt}
{\color{muted}Revenue and Pricing Manager | Pricing \textperiodcentered{} Capacity \textperiodcentered{} Demand Forecasting \textperiodcentered{} Revenue Assurance}\par\vspace{3pt}
{\small\color{muted}Cairo, Egypt \ifx\cvphone\empty\else\enspace\textperiodcentered\enspace \cvphone\ \fi\enspace\textperiodcentered\enspace \href{mailto:kareemfa5ry@gmail.com}{kareemfa5ry@gmail.com} \enspace\textperiodcentered\enspace \href{https://linkedin.com/in/kareem-fakhry}{linkedin.com/in/kareem-fakhry}}\par\vspace{2pt}
{\color{rule}\rule{\linewidth}{0.6pt}}

\section{Profile}
Data analyst with 15 years in airline revenue, pricing and operations. I build the full path from raw system files to a business decision: SQL and Python data pipelines, Power BI dashboards, forecasts, customer segmentation and statistical tests. I find money that is being lost, such as uncollected penalties, unbilled extra baggage and partners breaking contract rules, and I train commercial teams to use AI in their daily work.\par

\section{Skills}
\textbf{Revenue management:} Pricing strategy, capacity and inventory control, overbooking, demand forecasting, group and contract pricing\par\vspace{1pt}
\textbf{Revenue assurance:} Partner audits, penalty and fee checks, refund control, settlement with partners, audit procedures\par\vspace{1pt}
\textbf{SQL and databases:} PostgreSQL, MySQL, DBeaver, joins and window functions, materialized views, indexing and query tuning, triggers\par\vspace{1pt}
\textbf{BI and reporting:} Power BI, Power Query, Excel Data Model, star schema models, Excel and VBA, automated daily reports\par\vspace{1pt}
\textbf{Statistics and analysis:} Descriptive statistics, box plots and outlier checks, significance tests, customer segmentation, churn analysis\par\vspace{1pt}
\textbf{Source systems:} Amadeus (inventory, revenue accounting, billing files), Hitit, ATPCO, Accelya, SAP, Oracle\par\vspace{1pt}

\section{Experience}
\role{Data and AI Consultant (freelance)}{Nesma Airlines}{Cairo \textperiodcentered{} Feb 2026 -- Present}
\begin{itemize}
  \item Kept on as a consultant after the employee role ended. Built a log-processing pipeline, a daily forecasting pipeline and an audit application (see Projects).
  \item Since August 2026, running a paid AI programme for the commercial department: a 5-day, 15-hour bilingual Arabic and English foundation course, then advanced sessions on AI agents and automated workflows. Each participant builds and presents a working workflow from their own daily tasks.
\end{itemize}

\role{Revenue Manager}{Nesma Airlines}{Cairo \textperiodcentered{} Oct 2022 -- Jan 2026}
\rolenote{Joined after new owners bought the airline, with reporting systems switched off. Built the revenue and data team from one person to five.}
\begin{itemize}
  \item Restarted commercial operations in two to three months, rebuilding reporting by hand with six sales agents while systems were offline.
  \item Cleared a USD 1.5 million customer refund backlog before sales reopened. Settled about 60\% without cash, through future incentives and discounted peak-season seat blocks given to selected top agents.
  \item Turned one manual sheet into an automated daily report: capacity, passengers, utilisation, average fare and revenue, actual against plan, by route, sales channel, currency and booking window.
  \item Built a 136-field transaction table, one row per ticket action, combining four sales sources, with every money item kept in both original and reporting currency.
  \item Evaluated a commercial revenue management system (PROS), found it did not fit the Egyptian market, and built an in-house forecast instead.
  \item January 2023 to January 2024: seats doubled while revenue tripled from USD 1.6 million to USD 5.0 million, with average fare up 37\% and seat utilisation up from 71\% to 81\%.
\end{itemize}

\role{Head of Revenue and Network Planning}{Nile Air}{Cairo \textperiodcentered{} Jan 2017 -- Jun 2021}
\rolenote{Started as acting Revenue Manager and was confirmed after three months. Led four route controllers and yield executives.}
\begin{itemize}
  \item Set pricing strategy and repriced the full route network after the Egyptian currency devaluation.
  \item Brought fare filing in-house by securing a two-week ATPCO training programme, removing a paid outside dependency.
  \item Delivered interline agreements with flydubai, Pegasus and Hahn Air, and codeshare agreements with Air Cairo and EgyptAir, from contract to pricing setup and go-live.
  \item Took the Amadeus revenue accounting system through go-live with Finance and the vendor, fixing live data errors until reported figures could be trusted.
  \item Wrote the audit rules and procedures for a three-year outsourced revenue audit (Sutherland) covering sales, reissues, refunds and agent booking behaviour.
  \item Built the network and revenue data warehouse in SQL and Excel, and the five-year business plan prepared for the planned stock market listing.
  \item Decided with Planning which flights to cancel and which to operate through the COVID-19 closures and reopening.
\end{itemize}

\role{Air Arabia}{Egypt and UAE}{Dec 2010 -- Jan 2016}
\rolenote{Yield Management Executive (UAE, 2014--2016) \textperiodcentered{} Revenue Management Support Executive (UAE, 2014) \textperiodcentered{} Sales Support Executive (Egypt, 2012--2014) \textperiodcentered{} Call Centre Representative (Egypt, 2010--2012)}
\begin{itemize}
  \item Priced routes, controlled seat inventory and overbooking, and analysed no-shows and group demand. Handled the Egypt base across fares, sales and promotions.
  \item Supported the market study for the Jordan base launch, and handled all contracted Umrah group bookings across three airports for a full season.
\end{itemize}

\section{Projects}
\project{Competitor price tracker}{Nesma Airlines, 2023--2024}
\begin{itemize}
  \item Started as a script that captured fare calendar screenshots for the daily report. Grew into a scheduled scraper collecting one-way and 14-day return prices 90 days ahead for up to eight competing airlines per route.
  \item Cleaned data was emailed automatically every day to the commercial team, the revenue team and the chairman, and fed the pricing dashboards.
  \item The goal was daily market prices for every route, to support pricing decisions, inspired by earlier use of a paid fare-tracking service.
  \item Version 1 (2023) was a script that opened a browser, used a fare calendar view, took a screenshot per route, compressed the images and attached them to the daily report email for the commercial and revenue teams.
  \item Version 2 (late 2023 onward) was a full scraper deployed on the airline's own servers. Every day it collected prices for the target routes for the next three months, one-way and return with the return 14 days later.
  \item Output was a price table per route with one column per competing airline. Each raw record kept the reading date, sector, departure and return dates, carrier, airports, flight times and price.
  \item Each daily run sent a completion notice.
  \item Tools: Python, web scraping and Windows scheduled tasks.
\end{itemize}
\stack{Python, Web scraping, Scheduled tasks, Power Query}

\project{Daily revenue data pipeline and forecasting dashboard}{Nesma Airlines, 2026}
\begin{itemize}
  \item Loads daily ticket and flight exports into PostgreSQL and applies only the rows that changed, found by comparing checksums. Runs in 17 stages that restart from the point of failure; a normal daily run takes about 9 minutes.
  \item Forecasts every flight departing in the next 14 days by comparing its bookings so far with similar past flights, then scores each flight and flags the ones that need a pricing or inventory action.
  \item Moved the heavy processing from Python into the database after the first version needed 35 GB of memory and stalled.
  \item The problem was a daily answer for the Revenue Management team: which upcoming flights are filling too fast or too slowly.
  \item PostgreSQL does almost all the work; Python only runs the steps in order, times them, logs results and stops safely on failure. No pandas, Airflow or Spark.
  \item Each CSV export is loaded with PostgreSQL's COPY command into a landing table that is wiped and refilled in one transaction, so a failed load leaves the previous data untouched.
  \item The stages build calendar tables (including the Hijri calendar, which sets Ramadan, Hajj and Umrah demand), compare checksums to update only changed rows, and flag bad records instead of deleting them.
  \item Later stages refresh materialized views for exchange rates, flight totals, KPI summaries, yesterday's sales, sales by country and net position per ticket coupon.
  \item The final stage looks at each flight departing within 14 days. It builds the booking position by class and compares the booking curve with similar past flights at the same point before departure.
  \item It forecasts the final load, scores the flight and adds rule-based flags such as "high pickup" and "protect inventory".
  \item The dashboard is a React app: it asks a GraphQL server for data, the server checks a Redis cache and queries PostgreSQL if needed, and the page draws cards, charts and a map. Built with AI coding assistants.
\end{itemize}
\stack{PostgreSQL, Materialized views, Python, React, GraphQL, Redis}

\project{Customer and route analytics in Power BI}{Nesma Airlines}
\begin{itemize}
  \item Built dashboards for demand forecasting, customer segmentation, personalised offers and customer churn.
  \item Used averages, medians and spread measures, box plots to find outliers, and significance tests on sampled routes to check whether differences between routes were real before acting on them.
  \item The dashboards covered forecasting, customer segmentation, personalised offers and churn.
  \item Data was modelled as a star schema: a central fact table joined one-to-many to lookup tables, in Power BI and in Excel's Data Model, loaded with Power Query.
  \item Statistics were used to check before acting. Mean, median and spread described each route, box plots exposed outliers, and significance tests on sampled routes separated a real difference from chance, to prove or reject a hypothesis.
  \item Classification sorted customers into groups, which fed the segmentation and the personalised offers, and churn analysis showed who was likely to stop booking.
  \item The tools were Power BI, Power Query and Excel.
\end{itemize}
\stack{Power BI, Power Query, Star schema, SQL, Excel}

\project{Revenue audit engine}{Nesma Airlines}
\begin{itemize}
  \item Checked every ticket against filed fare rules and contracted fees. Flagged refunds without the required penalty, refunds on non-refundable fares, late cancellations and churning (an agent booking and cancelling the same seats again and again).
  \item Defined 38 standard breach categories and ran a manual check in parallel before relying on the automation.
  \item One quarter: 1,305 breaches across 170 sales partners, worth about USD 106,000, audited in weekly periods. About 957 were refunds, 171 voids and 100 commission cases. Finance accepted about 95\% of the cases submitted.
  \item The problem was finding agent rule breaches and raising debit memos at the airline's scale without an outside audit company.
  \item Two audit streams ran. The fare audit stored newly created fare bases and values in a table, then used rule logic and a lookup to flag differences between coded and used values.
  \item Example: promotional booking classes are non-refundable, so any ticket in them with a refund action is flagged.
  \item The billing audit checked GDS billing data for cost per segment per territory against contracted prices, plus churning and passive segments not cancelled.
  \item A Python desktop decoder reads the monthly GDS billing files and writes churning and passive-segment reports, a top-10 agents summary and a transaction report. Every decision is a fixed rule; no AI model is involved. Churning is repeated booking and cancelling beyond a set limit.
  \item Finance's roughly 95\% is an acceptance rate on submitted cases, not an accuracy measure; breaches the rules missed were not counted.
  \item Where cases were rejected, the usual cause was a rule that had not been properly filed, so there was nothing to charge against. The fix is upstream, in the rule data.
\end{itemize}
\stack{Excel \& VBA, Power Query, pandas, Amadeus billing files}

\project{Operational log parser and revenue checks}{Nesma Airlines, 2026}
\begin{itemize}
  \item Parsed 68 million raw airline operational log records covering 17 months, and linked 11 million ticket events across separate systems.
  \item Matched 99\% of 939,000 checked bags to the ticket that owned them, cutting unmatched records from 66\% to under 1\%. Found 40,981 paid extra-baggage records with no matching ticket, for revenue recovery.
  \item Confirmed 2,708 no-shows out of 3,779 cases checked, and set aside cases with conflicting boarding evidence so no passenger was wrongly charged.
  \item Made a key query about 660 times faster by adding a missing index. Protected results with locked tables, SHA-256 row fingerprints and 193 automated tests.
  \item The problem was raw airline operational logs that no one could trust. The aim was linked, trusted records, then lost revenue and wrong no-show charges found from them.
  \item The parser produced 11.17 million cross-system ticket links and 1.78 million coupon life-cycle timelines. All 51 of 51 reconciliation checks passed.
  \item Baggage: 931,769 of 938,632 checked bags matched to the owning ticket, and bag-tag matching precision rose from 67.2\% to 99.99\% once reused tags were handled.
  \item Missing ticket links fell from 65.7\% to 0.5\%. Check-in matched to flights in 99.81\% of cases across 7,886 flights.
  \item For 363,614 tickets with no boarding record, a four-level intended-flight finder raised coverage from 0\% to 69.1\%, with accuracy of 97.4\% to 99.9\% on known test cases.
  \item Joining the commercial data with airport data showed 3,225 coupons boarded on a different flight than ticketed.
  \item Controls: tables locked by state-machine triggers (created, running, completed), SHA-256 fingerprints of rows and files, one view per consumer, and 74 of 74 static checks passing alongside the 193 automated tests.
\end{itemize}
\stack{PostgreSQL, Triggers, Python, SHA-256, Automated tests}

\section{Education and training}
\begin{itemize}
  \item BSc Accounting, Alexandria University, Faculty of Commerce (English section), 2010
  \item ATPCO fare filing and coding, two-week instructor-led programme, Nile Air head office
  \item Business Process Outsourcing Skills Programme, Information Technology Institute (ITI), nine months
\end{itemize}

\section{Languages}
Arabic (native) \textperiodcentered{} English (full professional) \textperiodcentered{} French (basic)\par

\end{document}
